\documentclass{article}

\textwidth=6.5in
\textheight=8.5in
\voffset=-54pt
\oddsidemargin=0in
\evensidemargin=0in
\setlength{\parskip}{8pt}
\setlength{\parindent}{0pt}

\usepackage{amsmath, amssymb}
\usepackage{ifthen}

\usepackage[inline]{enumitem}
\setlist{topsep=0pt, parsep=8pt}

\usepackage[rgb,table]{xcolor}\convertcolorsUfalse
\usepackage{graphicx}

% change hyperref options
\PassOptionsToPackage{hyphens}{url}
\usepackage[bookmarks,colorlinks,linkcolor=black,citecolor=blue,pdfstartview=FitH,urlcolor=blue]{hyperref}
\hypersetup{pdfpagemode=UseNone}
\usepackage{bookmark}

\usepackage{graphicx}
\usepackage{multicol}
\usepackage{framed}
\usepackage{array}

\usepackage{icomma}

\usepackage{marginnote}
\usepackage[colorinlistoftodos]{todonotes}
\renewcommand{\marginpar}{\marginnote}

\usepackage{soul}

\usepackage{pgfplots}
\pgfplotsset{compat=1.18}  
\pgfplotscreateplotcyclelist{mycolor}{black, blue, red, green!70!black, magenta, purple, cyan, orange }  
\pgfplotsset{
  disabledatascaling,
  cycle list=tikzcolor, 
  axis lines=middle,
  every axis plot/.style={no marks, thick},
  every axis/.style={
    enlarge x limits=false,
    enlarge y limits=auto,
    xlabel style={at={(current axis.right of origin)},anchor=west},
    ylabel style={at={(current axis.above origin)},anchor=south},
    % extra x and y ticks for asymptotes
    extra x tick style={grid=major, grid style={dashed,black}},
    extra y tick style={grid=major, grid style={dashed,black}},
    /pgf/number format/1000 sep={},
  },    
  tick label style = {font=\footnotesize, fill=\currentbgcolor, inner sep=0pt},    
  xticklabel shift=1pt,    
  scaled ticks=false,
  scaled y ticks=false,
  unbounded coords=jump,
  samples=101,
  minor tick length=0,
  scatter/classes={
    open={mark=*,draw=black,fill=white, mark size=2, thin},
    closed={mark=*,draw=black,fill=black, mark size=2, thin}
  },
  width=3in,
}
\pgfplotsset{open/.style={only marks, mark=*, fill=white, thin, mark size=2}}
\pgfplotsset{closed/.style={only marks, mark=*, draw=black, fill=black,
    thin, mark size=2}}
\pgfplotsset{labeled/.style={nodes near coords, only marks, mark=*, draw=black, fill=black, thin, mark size=2, point meta=explicit symbolic}}

\usepackage{tikz}
\usetikzlibrary{
  matrix,%
  % enables the snake paths
  decorations.pathmorphing,%
  % enables bigger arrows
  decorations.markings,%
  decorations.pathreplacing,%
  decorations.text,%
  calc,%
  patterns,%
  shapes.geometric,%
  arrows,
  decorations.shapes,
  positioning,
  plotmarks,
  intersections
}
\tikzset{>=latex} % sets default arrow type in tikz
\tikzset{font=\normalsize}
\tikzset{mark size=1.5pt, mark options=thin}
\tikzset{pin distance=4pt,
  every pin edge/.style={<-, thin, shorten <= -2pt}}

\interfootnotelinepenalty=10000
\renewcommand{\thefootnote}{(\arabic{footnote})}

\usepackage{multirow, bigdelim}

% change to \usepackage[sols]{handout} to see solutions
\usepackage[sols, grading, notes]{handout}
\renewcommand{\workshopnote}{CA Notes.}    

\newcommand\iscurrentenumi[1]{%
  \ifthenelse{\equal{#1}{\theenumi}}{}{#1}}
\setenumerate[1]{ref=\#\arabic*}
\setenumerate[2]{ref=\protect\iscurrentenumi{\theenumi}(\alph*)}
\newcommand\enumiiprefix[2]{%
  \ifthenelse{{\equal{#1}{\theenumi}}\AND{\equal{#2}{(\alph{enumii})}}}{}{}%
  \ifthenelse{{\equal{#1}{\theenumi}}\AND{\NOT{\equal{#2}{(\alph{enumii})}}}}{#2}{}%
  \ifthenelse{\equal{#1}{\theenumi}}{}{#1#2}%
}
\setenumerate[3]{ref=\protect\enumiiprefix{\theenumi}{(\alph{enumii})}\roman*}

\makeatletter

% underline links               
\let\@href=\href
\renewcommand{\href}[2]{\@href{#1}{\ul{#2}}}

\makeatother

\definecolor{purple}{rgb}{0.5,0,1.0}

\newcommand{\psrnewpage}{\newpage}
\newcommand{\psreflection}[2][]{

  \ifthenelse{\not\equal{#1}{nocite}}{
    \begin{enumerate}[resume]
    \item
      \begin{problem}[1in]
        Please cite any help you got on this problem set:
        \begin{itemize}[nosep]
        \item If you worked with other students, please list their names here.
        \item If you used resources other than those on our Canvas site, please list them here.
        \end{itemize}
      \end{problem}

    \end{enumerate}
  }

  \probonly{%
    #2

    \psrnewpage

    \emph{Reflection space.} It's a good habit to take a few minutes at the end of each pset to reflect on what you learned and jot down notes to your future self. Here are a few reflection questions to get you started. Nothing you write here will be graded; it's just for you!
    \begin{itemize}
    \item Take a few minutes to look back at the problems. How could you explain to someone else how to do each problem? What was your strategy, and how did you know to use that strategy? If there are problems where you don't feel able to answer this, write down which ones they are. Come back to them in a few days when we post the homework solutions!

      \vspace{1.5in}

    \item Take a look back at the learning objectives on today's worksheet; how are you feeling about each one? If there are ones you know you need more practice with, write those down here.

      \vspace{1.5in}

    \item What did you find most challenging about this material? Do you have any lingering questions about it? If so, write them down here so you don't forget them. At some point, stop by office hours or MQC to ask!

      \vfill
    \end{itemize}      
  }
}

\usepackage{fancyhdr}
\lhead{\textcolor{white}{This content is protected and may not be shared, uploaded, or distributed.}}
\rhead{}
\rfoot{\vspace*{\fill}\footnotesize\textcopyright{} \emph{Harvard Math Ma}}
\renewcommand{\headrulewidth}{0pt}
\pagestyle{fancy}
\begin{document}

\begin{center}
  \large\textsc{Problem Set 21 - Finding Global Extrema}
\end{center}

There is no Edfinity for this problem set! As usual, be sure to justify your answers fully, and organize your work so the reader can follow it easily. 

\begin{enumerate}
\item Let $f(x) = x^3 - 6x^2 + 9x + 7$.

  \begin{enumerate}

  \item

    \begin{grading}
      2 points: 0.5 for finding $f'$, 1.5 for correctly solving for its zeros. 
    \end{grading}

    \begin{problem}[\vfill]
      Find the critical points of $f$. 
    \end{problem}

    \begin{solution}  % Jason Bland
      The critical points of $f$ are the $x$-values for which $f'(x)$ is
      zero or undefined.
      $f'(x) = 3x^2 - 12x + 9 = 3(x^2 - 4x + 3) = 3(x - 1)(x - 3)$, so the critical points are $\boxed{x = 1, 3}$. 
    \end{solution}

  \end{enumerate}

  \begin{grading}
    (b) and (c) together are 3 points: grade based on whether they have a correct strategy.
    \begin{itemize}[nosep]
    \item If they have correct answers but no strategy, give 0.5. 
    \item If they have a correct strategy but wrong answers because of arithmetic, deduct 0.25. 
    \end{itemize}
    If they use the First or Second Derivative Test, it's worth giving them the feedback that this is unnecessary.

    Other than the strategy in the solutions, they could also use the First Derivative Test or Second Derivative Test to classify the critical points. Then, they still need to compare $f(1)$ to $f(6)$ to decide where the global max is, and $f(3)$ to $f(-2)$ to decide where the global min is. Deduct 1.5 if they fail to do these checks. 
  \end{grading}

  \begin{enumerate}[resume]
  \item

    \begin{problem}[\stretch{1.25}]
      Does $f$ have a global maximum on $[-2, 6]$? If so, where (at what $x$-value)? If not, why not? Justify your answer. 
    \end{problem}

    \begin{solution}
      Since we're looking at a continuous function on a closed interval, $f$ must have a global maximum and global minimum on $[-2, 6]$. To find these global extrema, we can plug the critical points and endpoints into $f$:
      \begin{align*}
        f(-2) &= -43 \\
        f(1) &= 11 \\
        f(3) &= 7 \\
        f(6) &= 61
      \end{align*}
      So the \fbox{global maximum occurs at $x = 6$}.
    \end{solution}

  \item

    \begin{grading}
      1 point
    \end{grading}

    \begin{problem}[\nstretch{0.5}]
      Does $f$ have a global minimum on $[-2, 6]$? If so, where (at what $x$-value)? If not, why not? Justify your answer. 
    \end{problem}

    \begin{solution}
      Based on our work from the previous part, the \fbox{global minimum occurs at $x = -2$}.
    \end{solution}

  \end{enumerate}

\item  

  Let $g(t) = \dfrac{t^2}{t^4 + 16}$.

  \begin{enumerate}
  \item

    \begin{workshop}
      This is a good opportunity to have students think through whether they prefer the First Derivative Test or Second Derivative Test. (The Second Derivative Test seems annoying to me because calculating $g''$ is lengthy, but it's still totally valid to use it!) 
    \end{workshop}

    \begin{grading}
      \begin{itemize}[nosep]
      \item 1 for calculating $g'$ correctly (You'll see both the work for the Quotient Rule or Chain Rule below)
      \item 1.5 for solving $g' = 0$ correctly 
      \item 3 points for classifying the critical points using a correct method 
      \end{itemize}
    \end{grading}

    \begin{problem}[\stretch{3}]
      Find the critical points of $g$, and classify each (as a local minimum, local maximum, or neither). Make sure to explain how you come up with your classification. 
    \end{problem}

    \begin{solution}
      The critical points of $g$ are the points in the domain where $g'$ is
      undefined or zero. To find $g'(t)$, we use the Quotient Rule:
      \begin{align*}
        g'(t) &= \frac{(2t)(t^4+16)-(t^2)(4t^3)}{(t^4+16)^2} \\
        &= \frac{-2t^5+32t}{(t^4+16)^2} \\
        \intertext{Now we can simplify and factor.}
        &= \frac{-2t(t^4-16)}{(t^4+16)^2}\\
        &= \frac{-2t(t^2-4)(t^2+4)}{(t^4+16)^2} \\
        &= \frac{-2t(t-2)(t+2)(t^2+4)}{(t^4+16)^2}
      \end{align*}
      From here we can determine that $g'(t)$ is never undefined, and
      $g'(t)=0$ when $t=-2$, $0$, or $2$. So the critical points are
      \fbox{$-2$, $0$, and $2$}.\\

      It looks like it will be messy to find the $g''(t)$,
      so let's use the First Derivative Test to classify the critical points.
      Let's determine the sign of $f'(t)$ on the various intervals by testing points,
      noting that the denominator is always positive.
      \begin{align*}
        f'(-3) &= \frac{-2(\text{neg})(\text{neg})(\text{neg})(\text{pos})}{\text{positive}}
          = \text{positive} \\
        f'(-1) &= \frac{-2(\text{neg})(\text{neg})(\text{pos})(\text{pos})}{\text{positive}}
          = \text{negative} \\
        f'(1) &= \frac{-2(\text{pos})(\text{neg})(\text{pos})(\text{pos})}{\text{positive}}
          = \text{positive} \\
        f'(3) &= \frac{-2(\text{pos})(\text{pos})(\text{pos})(\text{pos})}{\text{positive}}
          = \text{negative} \\
      \end{align*}
      There is a sign chart below, but first we'll look at how to approach this problem using the Chain Rule...

      {\bf Here is what your work looks like if you use the Chain Rule on $g(t)$.} First, we'll rewrite the function as $g(t)=t^2\cdot (t^4+16)^{-1}$.
      \begin{align*}
          g'(t)&= 2t(t^4+16)^{-1}+t^2(-(t^4+16)^{-2}(4t^3)\\
          &=\frac{2t}{t^4+16}-\frac{4t^5}{(t^4+16)^2}\\
          &=\frac{2t}{t^4+16}\cdot\frac{t^4+16}{t^4+16}-\frac{4t^5}{(t^4+16)^2}\\
          &=\frac{2t(t^4+16)-4t^5}{(t^4+16)^2}\\
          &=\frac{2t^5+32t-4t^5}{(t^4+16)^2}\\
          &=\frac{32t-2t^5}{(t^4+16)^2}
      \end{align*}
This is a much more simplified numerator. The denominator is always positive, so changes of sign will be entirely controlled by the numerator (if they happen at all)
\begin{align*}
    0&=\frac{32t-2t^5}{(t^4+16)^2}\\
    0&= 32t-2t^5\\
    0&= 2t(16-t^4)
\end{align*}
We see that $g'(t)=0$ when $t=0,\;-2,\;2$. When $t<-2$, but factors, $2t$ and $16-t^4$ are negative, so the overall product is positive. When $-2<t<0$, $2t<0$ and $16-t^4>0$, so the overall product is negative. When $0<t<2$, both factors are positive, so the overall product is positive. And when $t>2$, $2t>0$ and $16-t^4<0$, so the overall product is negative.
      
      {\bf Regardless of what derivative rules you used in your approach, now we can draw a sign chart for $g'(t)$. }

      \begin{center}
        \begin{tikzpicture}[scale=0.8]
          \draw[<->] (-4,0) -- (4,0);
          \node[left] at (-4,0) {$f'(x)$};
          \foreach \x in {-2,0,2}
            \draw[shift={(\x,0)},color=black] (0pt,3pt) -- (0pt,-3pt) node[below] {$\x$};
          \foreach \x in {-2, 0, 2}
            \node[circle, fill=black, inner sep=2pt] at (\x,0) {};
          \node[above] at (-3,0) {$+$};
          \node[above] at (-1,0) {$-$};
          \node[above] at (1,0) {$+$};
          \node[above] at (3,0) {$-$};
        \end{tikzpicture}
      \end{center}
      Based on the sign chart, \fbox{$t=-2$ and $t=2$ are local maxima}, 
      and \fbox{$t=0$ is a local minimum}.
    \end{solution}

  \item

    \begin{workshop}
      It's fine to check the end behavior here, but it's also worth helping students understand why it isn't necessary: since $x = 2$ is the only critical point in $[1, \infty)$ and we know it's a local maximum, $g$ must be increasing on $(1, 2)$ and decreasing on $(2, \infty)$. So, 2 must be the global maximum on $[1, \infty)$.
    \end{workshop}

    \begin{grading}
      1.5 points: 0.5 for the answer, 1 for the justification. If they check end behavior, it's worth letting them know it isn't necessary and why. 
    \end{grading}

    \begin{problem}[\vfill\newpage]
      Does $g$ have a global maximum on $[1, \infty)$? If so, where? Justify your answer. 
    \end{problem}

    \begin{solution}
      Yes, and it occurs at $t=2$. We know this because there is only one critical
      point on the interval $[1, \infty)$ and the critical point is a local maximum,
      so it must be the global maximum on that interval.
    \end{solution}

  \item

    \begin{grading}
      2 points: 0.5 for evaluating $f(1)$, 1 for evaluating $\displaystyle \lim_{x \to \infty} f(x)$, 0.5 for drawing the correct conclusion 
    \end{grading}

    \begin{problem}[\vfill]
      Does $g$ have a global minimum on $[1, \infty)$? If so, where? Justify your answer. 
    \end{problem}

    \begin{solution}
      To determine whether there is a global minimum, we need to check the end behavior.
      At the left endpoint, $g(1)=\frac{1}{17}$. On the right side, we need to evaluate
      $\lim_{t \to \infty} g(t)$.
      \begin{align*}
        \lim_{t \to \infty} g(t) &= \lim_{t \to \infty} \frac{t^2}{t^4+16} \\
        \intertext{As $t\to\infty$, the dominant term in the denominator is $t^4$.}
        &= \lim_{t \to \infty} \frac{t^2}{t^4} \\
        &= \lim_{t \to \infty} \frac{1}{t^2} \\
        &= 0
      \end{align*}
      Based on this, we can say that $g$ does not have a global minimum on $[1, \infty)$.
    \end{solution}

  \end{enumerate}

  \begin{problemonly}
    \emph{Note:} You should have found 3 critical points. If you didn't, go back and fix it!
  \end{problemonly}

\item  

  \begin{grading}
    5 points: 1 for satisfying each condition.

    Since we didn't specify whether $f$ is differentiable on $[0, 10]$, they can have $f'(5)$ be 0 or undefined. 
  \end{grading}

  \begin{problem}[\nstretch{1.5}]
    Sketch the graph of a function $f$ on $[0, 10]$ satisfying all of the following properties.
    \begin{itemize}[nosep]
    \item $f$ is continuous on $[0, 10]$. 
    \item $f$ has two critical points on $(0, 10)$.
    \item $x = 5$ is a critical point of $f(x)$, but it is neither a local minimum nor a local maximum.
    \item The global maximum of $f$ on $[0, 10]$ is at $x = 3$.
    \item The global minimum of $f$ on $[0, 10]$ is at $x = 0$. 
    \end{itemize}
  \end{problem}

  \begin{solution}
    Here are two possibilities.
    \begin{center}
        \begin{tikzpicture}
          \begin{axis}[ytick=\empty, xtick={0,3,5,10}] 
            \addplot+[domain=0:4] {-5*(x-3)^2+45};
            \addplot[domain=4:5] {5*(x-5)^2+35};
            \addplot[domain=5:10] {-(x-5)^2+35};
            \addplot[closed] coordinates {(0,0) (10,10)};
          \end{axis}
        \end{tikzpicture}
        \hspace{0.5in}
        \begin{tikzpicture}
          \begin{axis}[ytick=\empty, xtick={0,3,5,10}] 
            \addplot+[domain=0:5] {-5*(x-3)^2+45};
            \addplot[domain=5:10] {-2*(x-5)+25};
            \addplot[closed] coordinates {(0,0) (10,15)};
          \end{axis}
        \end{tikzpicture}
      \end{center}
  \end{solution}

\item  

  Vishwesh is analyzing a function $f(x)$ that's continuous on $(-\infty, \infty)$. He's found that its critical points are $x = -3, -1, 2, 6$, and he's made the following sign chart for $f'$:
  \begin{center}
    \begin{tikzpicture}
      \begin{axis}[axis y line=none, x axis line style={-}, xtick={-3, -1, 2, 6}, ymin=-2, height=1in, width=4in, clip=false, tick label style={font=\small}]
        \addplot+[domain=-9:9, very thin]{0};
        \draw (-9, -1.5) node[right] {sign of $f'$};
        \draw (-4.5, -1.5) node{$-$};%
        \draw (-2, -1.5) node{$+$};%
        \draw (0.5, -1.5) node{$+$};%
        \draw (4, -1.5) node{$-$};%
        \draw (7.5, -1.5) node{$+$};%
      \end{axis}
    \end{tikzpicture}   
  \end{center}

  Answer each of the following questions with ``definitely yes'', ``definitely no'', or ``we don't have enough information''. If the answer is ``definitely yes'', say where the global max/min could be (list all possible $x$-values). If the answer is ``definitely no'' or ``we don't have enough information'', explain why. 

  \begin{enumerate}
  \item

    \begin{grading}
      1.5 points: 0.5 for knowing it does, 1 for saying where it could occur
    \end{grading}

    \begin{problem}[\vfill]
      Does $f$ have a global maximum on $[-3, 5]$? 
    \end{problem}

    \begin{solution}
      Definitely yes. On $[-3, 5]$, $f$ increases from $x=-3$ to $x=2$ and
      then decreases from $x=2$ to $x=5$, so the global maximum is at $x=2$. 
    \end{solution}

  \item

    \begin{grading}
      1.5 points: 0.5 for knowing it does, 1 for saying where it could occur
    \end{grading}

    \begin{problem}[\vfill]
      Does $f$ have a global minimum on $[-3, 5]$? 
    \end{problem}

    \begin{solution}
      Definitely yes. Since $f$ increases 
      from $x=-3$ to $x=2$ and then decreases from $x=2$ to $x=5$, the global minimum
      could occur at $x=-3$ or at $x=5$. 
    \end{solution}

  \item

    \begin{grading}
      1.5 points: 0.5 for answer, 1 for explanation 
    \end{grading}

    \begin{problem}[\vfill\newpage]
      Does $f$ have a global maximum on $(3, \infty)$?
    \end{problem}

    \begin{solution}
      Definitely no. On $(3, \infty)$, $f$ decreases from $x=3$ to $x=6$ and
      increases from $x=6$ onwards. Since there are no included endpoints, $f$
      will not have a global maximum.
    \end{solution}

  \item

    \begin{grading}
      1.5 points: 0.5 for knowing it does, 1 for knowing where it occurs
    \end{grading}

    \begin{problem}[\vfill]
      Does $f$ have a global minimum on $(3, \infty)$?
    \end{problem}

    \begin{solution}
      Definitely yes. On $(3, \infty)$, $f$ decreases from $x=3$ to $x=6$ and
      increases from $x=6$ onwards. The global minimum is at $x=6$.
    \end{solution}

  \item

    \begin{grading}
      2 points: 1 for knowing we don't have enough information, 1 for explaining why.

      If they think the answer is ``definitely no'' and say something about end behavior, give them 0.5 points (out of 2).
    \end{grading}

    \begin{problem}[\vfill]
      Does $f$ have a global maximum on $(-\infty, \infty)$?
    \end{problem}

    \begin{solution}
      We don't have enough information. It's possible that the global maximum could
      occur at $x=2$, but it's also possible that there is no global maximum, for example
      if $f$ increases without bound after $x=6$. 
    \end{solution}

  \item

    \begin{grading}
      1.5 points: 0.5 for knowing it does, 1 for saying where it could occur
    \end{grading}

    \begin{problem}[\vfill\newpage]
      Does $f$ have a global minimum on $(-\infty, \infty)$?
    \end{problem}

    \begin{solution}
      Definitely yes. The global minimum must occur at either $x=-3$ or at $x=6$. 
    \end{solution}

  \end{enumerate}

\item  

  A farmer has 240 feet of fence, which he will use to build two fenced-in pens, one shaped like a square and one shaped like a circle. The square pen is for his pet goat, and the sides need to be at least 10 feet long. The circular pen is for his pet llama, and it needs to have a radius of at least 8 feet.

  It would be natural to try to maximize the total area of the two pens; to do that, we first need to model the total area as a function of a single variable. That modeling is what you'll do in this problem. As usual, please follow the guidelines from \href{https://canvas.harvard.edu/files/20704944/download?download_frd=1}{Problem Set 6} in presenting your work. 

  \begin{enumerate}
  \item
    \begin{grading}
      5 points:
      \begin{itemize}[nosep]
      \item 1 for stating goal
      \item 1 for a labeled picture
      \item 1 for knowing that $A = s^2 + \pi r^2$ where $A$ is total area, $s$ is side of square, and $r$ is radius
      \item 1 for using the fact that there are 240 feet of fence to relate $s$ and $r$
      \item 1 for solving for $s$ in terms of $r$ and plugging that into the formula for $A$
      \end{itemize}
    \end{grading}

    \begin{problem}[\vfill]
      Express the total area of the two pens as a function of the radius of the circular pen.
    \end{problem}

    \begin{solution}
      The goal is to write down the total area of the two pens as a function of the 
      radius of the circular pen. We know that the farmer has $240$ feet of fence.

      Let's draw a picture and define variables.

      \begin{center}
        \begin{tikzpicture}[scale=2]
          \draw[very thick] (0,0) -- (1,0) -- (1,1) -- (0, 1) -- cycle;
          \node[left] at (0,0.5) {$s$};
          \draw[very thick]  (2,0.5) circle (0.5);
          \draw[fill=black] (2,0.5) circle (1pt);
          \draw (2,0.5) -- node[midway, above] {$r$} (2.5,0.5);
        \end{tikzpicture}
      \end{center}

      Let's call the radius of the circular pen $r$ and the side length of the
      square pen $s$. Let's first write down a formula for the total area. If
      we call the total area $A_\text{total}$, then
      \[A_\text{total} = A_\text{square} + A_\text{circle}\]
      where $A_\text{square}$ is the area of the square pen
      and $A_\text{circle}$ is the area of the circular pen. We know
      that $A_\text{square}=s^2$ and $A_\text{circle}=\pi r^2$, so we can
      substitute these in to our formula.
      \[A_\text{total} = s^2 + \pi r^2\]
      Now we need to express this as a function of just $r$. To do so, we need to
      express $s$ in terms of $r$. We can use the fact that the farmer has 240 feet
      of fence. The amount of fence he'll need for the square pen is $4s$, and the
      amount of fence he'll need for the circular pen is $2\pi r$.
      \begin{align*}
          4s + 2\pi r &= 240 \\
          4s &= 240 - 2\pi r \\
          s &= 60 - \tfrac{1}{2}\pi r
      \end{align*}
      Now that we have an expression for $s$ in terms of $r$, we can substitute
      this into our formula for the total area to get a function of $r$.
      \[A_\text{total} = (60 - \tfrac{1}{2}\pi r)^2 + \pi r^2\]
    \end{solution}

  \item

    \begin{grading}
      1.5 points: 0.5 for having 8 as the lower endpoint, 1 for having $\frac{100}{\pi}$ as the upper endpoint
    \end{grading}

    \begin{problem}[\nstretch{0.5}]
      What's the domain of the function you wrote in (a)? Explain. 
    \end{problem}

    \begin{solution}
      The domain will be the possible values of $r$, the radius of the circular pen.
      We're told in the problem that the radius must be at least $8$ feet, so we just
      need to determine the maximum possible value of the radius.
      
      We can't
      make the radius arbitrarily large, because we only have $240$ feet of fence.
      If we use a lot of fence for the circular pen, we won't have very much for the
      square pen. Thus, the maximum possible radius is determined by the minimum possible
      side length of the square pen, $10$ feet. We can use the relationship we found between $s$ and $r$ to solve for $r$ when $s=10$. 
      \begin{align*}
          4(10) + 2\pi r &= 240 \\
          2\pi r &= 200 \\
          r &= \frac{100}{\pi}
      \end{align*}
      This shows that the maximum possible radius is $\frac{100}{\pi}$.

      Therefore the domain is $8\leq r \leq \frac{100}{\pi}$.
    \end{solution}

  \end{enumerate}

\item % 6.4.9, modified

  Amelia is a potter who makes and sells bowls. If she prices her bowls at $\$10$ per bowl, she can sell 70 bowls every week. For each dollar she increases the price, the number of bowls she sells in a week decreases by 5. 

  \begin{enumerate}

  \item

    \begin{grading}
      2 points
    \end{grading}

    \begin{problem}[\stretch{0.75}]
      Express the number of bowls Amelia sells each week as a function of the price she charges per bowl. 
    \end{problem}

    \begin{solution}
      There is a constant rate of change,
      so we know that this is a linear function. The slope of the graph is $-5$, and it
      passes through the point $(10, 70)$, so we can write down an equation
      using point-slope form.
      \[y - 70= -5(x-10) \]
      If $N$ is the number of bowls sold in a week and $p$ is the price per bowl 
      (in dollars per bowl), then
        \[N = -5(p-10)+70
        = -5p + 120.\]
    \end{solution}

  \item \label{6.4.9-revenue}

    \begin{grading}
      1 point
    \end{grading}

    \begin{problem}[\stretch{0.75}]
      Express Amelia's weekly revenue as a function of the price she charges per bowl. 
    \end{problem}

    \begin{solution} 
      Revenue is (price per bowl)(number of bowls sold), so if call the
      revenue in dollars $R$,
      \[R=pN=p(-5p+120)=-5p^2+120p.\]
    \end{solution}

  \item

    \begin{grading}
      3 points: 1 for finding the critical point, 2 for justifying why it's the global maximum. They can either check the endpoints, or since there's only one critical point, they can use the First Derivative Test or Second Derivative Test.

      To have a complete argument in the latter case, they should technically say that this works because there's only one critical point (and it ends up being a local maximum). If they don't say this, give them some feedback along with a 0.01 deduction (just so they see it). 
    \end{grading}

    \begin{problem}[\vfill]
      Amelia has decided that she will charge at least \$3 and at most \$20 per bowl. Where is the global maximum of the function you found in \ref{6.4.9-revenue} on $[3, 20]$? Justify your answer carefully.
    \end{problem}

    \begin{solution}
      We determined in part (b) that $R=f(p)$, where $f(p) = -5p^2+120p$. To find
      the global maximum of $f$ on $[3, 20]$, we need to first find the critical points.
      We know that $f'(p)=-10p+120$, and so the only critical point is $p=12$. When
      $p=12$, $f(12)=720$.

      We need to compare this value to the values of $f$ at the endpoints, $p=3$ and
      $p=20$. We can calculate that $f(3)=315$ and $f(20) = 400$. Based on this,
      the global maximum occurs at $p=12$.

      We also know that this must be the global maximum because $p=12$ is the only
      critical point on the interval $[3, 20]$, and since $f''(p)=-10$, it is a
      local maximum by the Second Derivative Test.
    \end{solution}

  \item

    \begin{grading}
      1 point. They get full credit if this is consistent with their answer to (c). 
    \end{grading}

    \begin{problem}[\newpage]
      Summarize your findings by filling in the blanks: if Amelia wants to maximize her weekly revenue, she should charge \$\fitb{0.5in}{12} per bowl. If she does this, her weekly revenue will be \$\fitb{1in}{720}. 
    \end{problem}

\item 
\begin{grading} 1 point.  
\end{grading}
\begin{problem}[\vfill] For this specific example, it's worth noting that we don't actually \textbf{need} calculus to find the maximum of the revenue function. Given that the revenue function you built is a quadratic, confirm the location of the vertex without using calculus (and only algebra skills)!\\
\end{problem}
\begin{solution}
    Our revenue function was written as $R=p(-5p+120)$, which allows us to find its zeros relatively quickly. These occur at $p=0$ and $p=24$. The vertex is positioned exactly in the center of these coordinates, at $p=12$. We can evaluate the revenue there to find how much money Amelia's maximum weekly revenue is: $R(12)=12(-5\cdot 12 +120) = 720$
\end{solution}
  \end{enumerate}

\end{enumerate}
\renewcommand{\psrnewpage}{{\centering\rule{5in}{0.4pt}}}

\psreflection{For next class, read \href{https://canvas.harvard.edu/files/20608914/download?download_frd=1}{\S 4.3 of \emph{Calculus} by Hughes-Hallett et. al} through Example 4 but skipping Example 2.}
\end{document}
